\documentclass[10pt,a4paper]{amsart}
\usepackage[margin=19mm]{geometry}
\usepackage{amsmath,amssymb,mathtools}
\usepackage[scr=rsfs,cal=euler]{mathalfa}
\usepackage{xfrac}
\usepackage[unicode,hidelinks,bookmarks=false]{hyperref}
\newcommand{\ie}{i.e.\ }
\newcommand{\cf}{cf.\ }
\newcommand{\resp}{respectively\ }
\newcommand{\Subsection}{\subsection}
\providecommand{\nobreakdash}{\nobreak\hbox{-}}
\setlength{\emergencystretch}{2em}
\multlinegap0pt
\usepackage{mathtools}
\usepackage{enumitem}
\newtheorem{theorem}{Theorem}[section]
\newtheorem{proposition}[theorem]{Proposition}
\newtheorem{lemma}[theorem]{Lemma}
\newtheorem{corollary}[theorem]{Corollary}
\newtheorem{definition}[theorem]{Definition}
\newtheorem{example}[theorem]{Example}
\newtheorem{remark}[theorem]{Remark}
\numberwithin{equation}{section}
\begin{document}
\title[Exchangeable Testing Against an Unknown Benchmark]{Exchangeable Testing Against an Unknown Benchmark}
\author{Alexander Gnedin}
\date{}
\maketitle
\noindent\emph{Source and licence.} Exchangeable Testing Against an Unknown Benchmark. Version arXiv:2608.04838v1 (CC BY 4.0). This material is distributed under the Creative Commons Attribution 4.0 International licence, http://creativecommons.org/licenses/by/4.0/, as recorded in the arXiv OAI licence metadata for this exact version and shown on the abstract page https://arxiv.org/abs/2608.04838v1. Full rights notice: licenses/arxiv-2608.04838-CC-BY-4.0.txt.\par\smallskip
\noindent\textit{Editorial scope.} Counted unit: the original \texttt{theorem} environment \texttt{-} at source lines 238--240 together with its complete proof at lines 242--244; the delivered document keeps source lines 191--245 so that every referenced label and every citation resolves inside it. Native editable source is the paper's own concatenated TeX; the AMS wrapper is editorial formatting only. No publisher class, upstream script or unrelated artwork is redistributed.
\begin{equation}\label{eq:LLN}
\frac{R_n}{n}\to X\quad\text{a.s.\ as }n\to\infty.
\end{equation}

\section{Identifiability}

This section is devoted to the analysis of parameters of the binary testing model. The parameters in our model are the size \(k\) of the initial players group and the prior distribution of winner's rank within the group. In this section we denote this discrete prior \((p_1,\dots,p_k)\); the set of such priors is a \(k\)-simplex (of dimension \(k-1\)), denoted \(\mathcal{P}_k\).

It is sometimes convenient to use parallel homogeneous coordinates, passing to odds signature, that is nonzero nonnegative sequences of weights \((w_1,\dots,w_k)\), which normalise by \(W:=\sum_{r=1}^k w_r\). The two pieces of notation are matched as
\[
(w_1,\dots,w_k)\propto(p_1,\dots,p_k),
\]
so we do not distinguish proportional odds signatures.

\begin{definition}
Two rank priors \((p_1,\dots,p_k)\) and \((q_1,\dots,q_\ell)\) are said to be binary-indistinguishable if they induce the (infinite) binary testing procedures with the same distribution. The definition extends to the representation of priors via their odds signatures.
\end{definition}

The partition of priors in indistinguishability classes of this equivalence relation is trivial within simplex \(\mathcal{P}_k\) of a given degree. Indeed, the Beta densities \(h_{k,r}\) of uniform order statistics comprise a positive Bernstein basis of the space of homogeneous polynomials of degree \(k-1\). Thus the set of such mixtures is a \((k-1)\)-dimensional simplex spanned on the extreme points \((h_{k,r},r\in[k])\). By de Finetti's theorem \eqref{eq:LLN}, the binary experiment identifies its continuous prior, hence the binary indistinguishability at fixed level amounts to equality \((p_1,\dots,p_k)=(q_1,\dots,q_k)\).

\begin{definition}
The elements of the simplex \(\mathcal{H}_k\), that is mixtures of Beta distributions \((h_{k,r},r\in[k])\), are called choice priors of degree \(k\).
\end{definition}

The mixing operation for fixed \(k\) is an isomorphism between \(\mathcal{P}_k\) and \(\mathcal{H}_k\), which allowed us on the latent level to distinguish the discrete priors. Now, the simplices of choice priors comprise an inductive system, \(\mathcal{H}_1\subset\mathcal{H}_2\subset\cdots\), where \(\mathcal{H}_1\) has a sole point \(h_{1,1}\) which is the uniform distribution. Looking across degrees \(k\) brings forward a new phenomenon: the discrete uniform priors, with signatures \((1),(1,1),(1,1,1),\dots\), all correspond to the continuous uniform distribution \(\mathrm{U}[0,1]\). Indeed, choosing at random one of the order statistics yields the uniform distribution, hence the basic Markov--P\'olya urn. Thus, geometrically, the barycenters of the simplices \(\mathcal{H}_k\) all coincide with \(h_{1,1}\), while the barycenters of \(\mathcal{P}_k\)'s live in different dimensions despite producing the same binary processes. Probabilistically, this mismatch is easy to explain: testing against a player randomly chosen from \(P_{-1},\dots,P_{-k}\) makes them exchangeable; a phenomenon valid also for finite series of tests, as long as we stay within the paradigm of comparing only with the winner. A more insightful for our study explanation is that the discrete uniform prior \(p_0(r)\) on \(\mathcal{P}_r\) through the ranking rule elevates to the uniform posteriors; in particular if we take a singleton prior group and update the rank of the benchmark through binary comparisons, the observed process will be precisely the same as the basic Markov--P\'olya urn or Kingman's paintbox with uniform split.

Invoking de Finetti's theorem again, the binary indistinguishability of \((w_1,\dots,w_k)\) and \((v_1,\dots,v_\ell)\), \(k<\ell\), amounts to the equality of expansions in the Bernstein bases
\begin{equation}\label{eq:equal}
\sum_{i=1}^k w_i\,h_{k,i}(x)=\sum_{j=1}^\ell v_j\,h_{\ell,j}(x).
\end{equation}
This relation has interpretation in terms of the general Kerov--Vershik theory of multiplicative branching graphs. The instance \cite{VershikKerov1987} in focus here is the Pascal graph canonically associated with the algebra of symmetric functions in two formal variables \(a,b\) via the Hausdorff degree elevation identity \((a+b)a^Ab^B=a^{A+1}b^B+a^Ab^{B+1}\) factored through \(a+b=1\).

Specialising \(a=x\), \(b=(1-x)\) as real numbers leads to the Bernstein basis conversion identity
\begin{equation}\label{eq:Bernstein}
h_{k,i}(x)=\sum_{j=i}^{\ell-k+i}\frac{\binom{j-1}{i-1}\binom{\ell-j}{k-i}}{\binom{\ell}{k}}h_{\ell,j}(x),
\end{equation}
where the hypergeometric coefficients come from the path-counting in the graph. By substituting this identity in the indistinguishability criterion \eqref{eq:equal} on the linearly independent higher-degree basis, we obtain the following combinatorial classification. Two mixtures are identical if and only if the odds signature of the higher degree is the hypergeometric projection of the lower-degree \(w\):
\begin{equation}\label{eq:vj}
v_j=\sum_{i=1}^k w_i\cdot\frac{\binom{j-1}{i-1}\binom{\ell-j}{k-i}}{\binom{\ell}{k}},\qquad\text{for }j=1,\dots,\ell.
\end{equation}

To translate the elevation technique in the language of testing against the benchmark, we introduce the operation of initial group extension.

\begin{definition}
Let \(P_{-k},\dots,P_{-1}\) be a group of players generating some benchmark rank prior \(p_0\). For integer \(\ell>k\), we say that \(P_{-\ell},\dots,P_{-k},\dots,P_{-1}\) is a prior group extension if the players \(P_{-\ell},\dots,P_{-k-1}\) joined the primary group successively, and the distribution of the winner's rank in the whole growing group was updated according to the rank rule. The resulting benchmark rank distribution \(q_0\) on \([\ell]\) is called extension of the prior \(p_0\).
\end{definition}


\begin{theorem}
For \(\ell>k\) two rank priors \(p_0,q_0\) of degree \(k\) and \(\ell\), respectively, are binary-indistinguishable if and only if \(q_0\) is the extension of \(p_0\).
\end{theorem}

\begin{proof}
The Bernstein basis conversion formula shows that the continuous density generated by \(p_0\) of degree \(k\) expands uniquely in the higher-degree basis of degree \(\ell\) with coefficients given by \eqref{eq:vj}. By definition, the extension \(q_0\) is the law of the running rank after the additional \(\ell-k\) players have been inserted according to the ranking rule; the path-counting argument that produces \eqref{eq:vj} is therefore identical to the combinatorial construction of the extension. Hence the two continuous densities coincide if and only if \(q_0\) is that extension.
\end{proof}


\begin{thebibliography}{99}
\bibitem[VershikKerov1987]{VershikKerov1987} A.~M.~Vershik and S.~V.~Kerov, Locally semisimple algebras. Combinatorial theory and the \(K_0\)-functor, \textit{Journal of Soviet Mathematics} \textbf{38} (1987) 1701--1733.
\end{thebibliography}
\end{document}
